\begin{table}[!t]
\centering\footnotesize
\setlength{\tabcolsep}{4pt}
\caption{Uç Nokta Ablasyonları: Varyant ile Tam Yöntem Arasındaki Son Doğruluk Farkı (Yüzde Puan)}\label{tab:ablation}
\begin{adjustbox}{max width=\columnwidth}
\begin{tabular}{lrrrr}
\toprule
Varyant & Hücre & Ort. & En az & En çok \\
\midrule
Fed-MDBSCAN-G, yalnız 1. aşama & 12 & $-$0.081 & $-$1.50 & $+$0.84 \\
Fed-MDBSCAN-G, etkinleşme belleği yok & 12 & 0.000 & 0.00 & 0.00 \\
Fed-MDBSCAN-G, valf yok & 12 & $+$0.278 & $-$0.27 & $+$3.03 \\
\bottomrule
\end{tabular}
\end{adjustbox}

\par\vspace{3pt}\parbox{\linewidth}{\footnotesize Varyant eksi tam bileşik yöntem, yüzde puan cinsinden son doğruluk farkı; hücreler eşleşen ablasyon bloğunda üç tohum üzerinden veri kümesi~$\times$~koşul ortalamalarıdır. Aralıklar koşullar arasındaki değişkenliği gösterir, güven aralığı değildir. Ek belge, eşleşmiş yarıçaplı bir kontrolü ve Aşama-1 yarıçap çarpanına duyarlılığı ekler.}
\end{table}
